\documentclass[10pt]{extarticle}
\usepackage[a4paper,margin=1.5cm]{geometry}
\usepackage[T1]{fontenc}
\usepackage{amsmath,amssymb}
\usepackage{graphicx}
\usepackage[export]{adjustbox}
\usepackage{multicol}
\usepackage{enumitem}
\usepackage{booktabs,array,tabularx}
\setlength{\columnsep}{0.7cm}
\setlength{\parindent}{0pt}
\setlength{\parskip}{3pt}
\pagestyle{empty}
\makeatletter
\newcommand{\qfigdef}[4]{\@namedef{qf@file@#1}{#2}\@namedef{qf@w@#1}{#3}\@namedef{qf@wide@#1}{#4}}
\newcommand{\qf@check}[1]{\@ifundefined{qf@file@#1}{\PackageError{score2}{Unknown figure id #1}{}}{}}
\newcommand{\qf@img}[4][]{\includegraphics[width=#2\dimexpr\@nameuse{qf@w@#4}pt\relax,max width=#3,max height=0.45\textheight,#1]{\@nameuse{qf@file@#4}}}
\newif\ifqf@full
\newif\ifqf@boxed
\newsavebox{\qf@box}
\newcommand{\qf@open}{\global\qf@fulltrue\par\global\qf@boxedtrue
  \setbox\qf@box\vbox\bgroup\hsize\textwidth\linewidth\textwidth}
\newcommand{\qf@emit}[1]{\par\egroup\global\qf@boxedfalse
  \ifdim\dimexpr\ht\qf@box+\dp\qf@box\relax<#1\textheight\noindent\usebox{\qf@box}\par\else\unvbox\qf@box\fi}
\newcommand{\qfullsplit}{\ifqf@boxed\qf@emit{0.85}\fi}
\newcommand{\qf@close}{\ifqf@boxed\qf@emit{0.45}\fi\global\qf@fullfalse}
\newcommand{\qfullbegin}{\end{multicols}\qf@open}
\newcommand{\qfullend}{\qf@close\begin{multicols}{2}}
\newcommand{\qfullstart}{\qf@open}
\newcommand{\qfulldone}{\qf@close}
\newcommand{\qfig}[2][1]{\qf@check{#2}%
  \ifqf@full\qf@blk{#1}{#2}\else
  \ifnum\@nameuse{qf@wide@#2}=1 \expandafter\@firstoftwo\else\expandafter\@secondoftwo\fi
  {\end{multicols}\begin{center}\qf@img{#1}{\textwidth}{#2}\end{center}\begin{multicols}{2}}%
  {\qf@blk{#1}{#2}}\fi}
\newcommand{\qf@blk}[2]{\par\smallskip{\centering\qf@img{#1}{\linewidth}{#2}\par}\smallskip}
\newcommand{\qfigi}[2][1]{\qf@check{#2}\qf@img[valign=c]{#1}{\linewidth}{#2}}
\makeatother
\newcommand{\papersec}[1]{\par\bigskip{\centering\Large\bfseries PAPER-#1\par}\medskip}
\newcommand{\partsec}[2]{\par\medskip\noindent{\large\bfseries PART-#1: #2}\par\smallskip}
\newcommand{\sect}[1]{\par\smallskip\noindent\textbf{#1}\par}
\newcommand{\qn}[1]{\par\smallskip\noindent\textbf{#1.}~}
\newcommand{\paraq}[2]{\par\smallskip\noindent\textbf{Paragraph for Questions #1 and #2}\par}
\newenvironment{qoptions}{\begin{itemize}[leftmargin=1.8em,itemsep=1pt,topsep=2pt,label={}]}{\end{itemize}}
\newcommand{\opt}[2]{\item[(#1)] #2}
% figure registry (generated by score2_tools.py finalize)
\qfigdef{p1_phy_s1_q01_1}{JEE Advanced (2021 Series) - Test 10 (13 May 2021)-img/p1_phy_s1_q01_1.png}{105.3}{0}
\qfigdef{p1_phy_s1_q02_1}{JEE Advanced (2021 Series) - Test 10 (13 May 2021)-img/p1_phy_s1_q02_1.png}{107.0}{0}
\qfigdef{p1_phy_s1_q04_1}{JEE Advanced (2021 Series) - Test 10 (13 May 2021)-img/p1_phy_s1_q04_1.png}{56.6}{0}
\qfigdef{p1_phy_s1_q05_1}{JEE Advanced (2021 Series) - Test 10 (13 May 2021)-img/p1_phy_s1_q05_1.png}{106.3}{0}
\qfigdef{p1_phy_s2_q02_2}{JEE Advanced (2021 Series) - Test 10 (13 May 2021)-img/p1_phy_s2_q02_2.png}{103.5}{0}
\qfigdef{p1_phy_s2_q03_1}{JEE Advanced (2021 Series) - Test 10 (13 May 2021)-img/p1_phy_s2_q03_1.png}{99.6}{0}
\qfigdef{p1_phy_s2_q04_1}{JEE Advanced (2021 Series) - Test 10 (13 May 2021)-img/p1_phy_s2_q04_1.png}{85.2}{0}
\qfigdef{p1_phy_s2_q04_2}{JEE Advanced (2021 Series) - Test 10 (13 May 2021)-img/p1_phy_s2_q04_2.png}{89.1}{0}
\qfigdef{p1_phy_s2_q04_3}{JEE Advanced (2021 Series) - Test 10 (13 May 2021)-img/p1_phy_s2_q04_3.png}{88.4}{0}
\qfigdef{p1_phy_s2_q04_4}{JEE Advanced (2021 Series) - Test 10 (13 May 2021)-img/p1_phy_s2_q04_4.png}{85.2}{0}
\qfigdef{p1_phy_s3_q05_1}{JEE Advanced (2021 Series) - Test 10 (13 May 2021)-img/p1_phy_s3_q05_1.png}{101.1}{0}
\qfigdef{p1_phy_s3_q06_1}{JEE Advanced (2021 Series) - Test 10 (13 May 2021)-img/p1_phy_s3_q06_1.png}{130.8}{0}
\qfigdef{p1_phy_s3_q08_1}{JEE Advanced (2021 Series) - Test 10 (13 May 2021)-img/p1_phy_s3_q08_1.png}{158.0}{0}
\qfigdef{p1_chem_s1_q02_1}{JEE Advanced (2021 Series) - Test 10 (13 May 2021)-img/p1_chem_s1_q02_1.png}{291.1}{1}
\qfigdef{p1_chem_s1_q05_1}{JEE Advanced (2021 Series) - Test 10 (13 May 2021)-img/p1_chem_s1_q05_1.png}{29.4}{0}
\qfigdef{p1_chem_s1_q05_A}{JEE Advanced (2021 Series) - Test 10 (13 May 2021)-img/p1_chem_s1_q05_A.png}{32.6}{0}
\qfigdef{p1_chem_s1_q05_B}{JEE Advanced (2021 Series) - Test 10 (13 May 2021)-img/p1_chem_s1_q05_B.png}{31.4}{0}
\qfigdef{p1_chem_s1_q05_C}{JEE Advanced (2021 Series) - Test 10 (13 May 2021)-img/p1_chem_s1_q05_C.png}{33.0}{0}
\qfigdef{p1_chem_s1_q05_D}{JEE Advanced (2021 Series) - Test 10 (13 May 2021)-img/p1_chem_s1_q05_D.png}{54.2}{0}
\qfigdef{p1_chem_s2_q02_1}{JEE Advanced (2021 Series) - Test 10 (13 May 2021)-img/p1_chem_s2_q02_1.png}{99.9}{0}
\qfigdef{p1_chem_s2_q02_2}{JEE Advanced (2021 Series) - Test 10 (13 May 2021)-img/p1_chem_s2_q02_2.png}{99.4}{0}
\qfigdef{p1_chem_s2_q02_3}{JEE Advanced (2021 Series) - Test 10 (13 May 2021)-img/p1_chem_s2_q02_3.png}{95.0}{0}
\qfigdef{p1_chem_s2_q02_4}{JEE Advanced (2021 Series) - Test 10 (13 May 2021)-img/p1_chem_s2_q02_4.png}{98.4}{0}
\qfigdef{p1_chem_s3_q04_1}{JEE Advanced (2021 Series) - Test 10 (13 May 2021)-img/p1_chem_s3_q04_1.png}{187.6}{0}
\qfigdef{p1_chem_s3_q05_1}{JEE Advanced (2021 Series) - Test 10 (13 May 2021)-img/p1_chem_s3_q05_1.png}{307.5}{1}
\qfigdef{p1_chem_s3_q06_1}{JEE Advanced (2021 Series) - Test 10 (13 May 2021)-img/p1_chem_s3_q06_1.png}{85.7}{0}
\qfigdef{p2_phy_s1_q02_1}{JEE Advanced (2021 Series) - Test 10 (13 May 2021)-img/p2_phy_s1_q02_1.png}{102.1}{0}
\qfigdef{p2_phy_s1_q04_1}{JEE Advanced (2021 Series) - Test 10 (13 May 2021)-img/p2_phy_s1_q04_1.png}{53.6}{0}
\qfigdef{p2_phy_s1_q07_1}{JEE Advanced (2021 Series) - Test 10 (13 May 2021)-img/p2_phy_s1_q07_1.png}{128.9}{0}
\qfigdef{p2_phy_s2_para1_2_1}{JEE Advanced (2021 Series) - Test 10 (13 May 2021)-img/p2_phy_s2_para1_2_1.png}{130.3}{0}
\qfigdef{p2_phy_s3_q01_1}{JEE Advanced (2021 Series) - Test 10 (13 May 2021)-img/p2_phy_s3_q01_1.png}{88.7}{0}
\qfigdef{p2_phy_s3_q02_1}{JEE Advanced (2021 Series) - Test 10 (13 May 2021)-img/p2_phy_s3_q02_1.png}{241.0}{0}
\qfigdef{p2_phy_s3_q04_1}{JEE Advanced (2021 Series) - Test 10 (13 May 2021)-img/p2_phy_s3_q04_1.png}{83.2}{0}
\qfigdef{p2_phy_s3_q05_1}{JEE Advanced (2021 Series) - Test 10 (13 May 2021)-img/p2_phy_s3_q05_1.png}{157.3}{0}
\qfigdef{p2_phy_s3_q07_1}{JEE Advanced (2021 Series) - Test 10 (13 May 2021)-img/p2_phy_s3_q07_1.png}{95.1}{0}
\qfigdef{p2_phy_s3_q08_1}{JEE Advanced (2021 Series) - Test 10 (13 May 2021)-img/p2_phy_s3_q08_1.png}{91.1}{0}
\qfigdef{p2_chem_s1_q03_1}{JEE Advanced (2021 Series) - Test 10 (13 May 2021)-img/p2_chem_s1_q03_1.png}{262.5}{1}
\qfigdef{p2_chem_s1_q04_A}{JEE Advanced (2021 Series) - Test 10 (13 May 2021)-img/p2_chem_s1_q04_A.png}{60.5}{0}
\qfigdef{p2_chem_s1_q04_B}{JEE Advanced (2021 Series) - Test 10 (13 May 2021)-img/p2_chem_s1_q04_B.png}{102.1}{0}
\qfigdef{p2_chem_s1_q04_C}{JEE Advanced (2021 Series) - Test 10 (13 May 2021)-img/p2_chem_s1_q04_C.png}{107.0}{0}
\qfigdef{p2_chem_s1_q04_D}{JEE Advanced (2021 Series) - Test 10 (13 May 2021)-img/p2_chem_s1_q04_D.png}{100.7}{0}
\qfigdef{p2_chem_s1_q05_1}{JEE Advanced (2021 Series) - Test 10 (13 May 2021)-img/p2_chem_s1_q05_1.png}{309.3}{1}
\qfigdef{p2_chem_s1_q05_2}{JEE Advanced (2021 Series) - Test 10 (13 May 2021)-img/p2_chem_s1_q05_2.png}{62.0}{0}
\qfigdef{p2_chem_s1_q05_3}{JEE Advanced (2021 Series) - Test 10 (13 May 2021)-img/p2_chem_s1_q05_3.png}{87.2}{0}
\qfigdef{p2_chem_s3_q05_1}{JEE Advanced (2021 Series) - Test 10 (13 May 2021)-img/p2_chem_s3_q05_1.png}{169.2}{0}
\qfigdef{p2_chem_s3_q06_1}{JEE Advanced (2021 Series) - Test 10 (13 May 2021)-img/p2_chem_s3_q06_1.png}{195.7}{0}

\begin{document}
\begin{center}{\Large\bfseries JEE Advanced (2021 Series) --- Test 10 (13 May 2021)}\end{center}
\begin{multicols}{2}
\papersec{1}
\partsec{1}{PHYSICS}
\sect{SEC--1}
\qn{1} Two vertical strings having different lengths are joined to a uniform rod at points $A$ and $B$. The rod is at rest. Now string connected to point $B$ is burnt. Which of the following is/are correct just after this ?
\qfig{p1_phy_s1_q01_1}
\begin{qoptions}
\opt{A}{Acceleration of COM is horizontal.}
\opt{B}{Acceleration of COM is vertical.}
\opt{C}{Acceleration of particle at $A$ is zero}
\opt{D}{Acceleration of particle at $A$ is horizontal}
\end{qoptions}
\qn{2} A simple pendulum of length $\ell=30$ cm and mass $m=2$kg is attached to a cart of mass $3$ kg. The cart can move frictionlessly on level ground. The bob is released from rest when string makes an angle of $1.8^\circ$ with the horizontal :-
\qfig{p1_phy_s1_q02_1}
\begin{qoptions}
\opt{A}{The amplitude of oscillation of cart is $\dfrac\pi5$ cm.}
\opt{B}{The time period of oscillation is $\dfrac\pi5$ sec}
\opt{C}{The cart and bob oscillate with phase difference of $\pi$.}
\opt{D}{The time period of oscillation is $\dfrac{\pi\sqrt3}{5}$ sec.}
\end{qoptions}
\qn{3} Consider the potential energy given by $U=\dfrac{x^4}{4}-\dfrac{x^3}{3}-3x^2$
\begin{qoptions}
\opt{A}{$x=0$ is position of stable equilibrium.}
\opt{B}{$x=3$ is position of stable equilibrium.}
\opt{C}{If body starts from $x=-2$, it can go to $x=3$ if its kinetic energy is greater than $\dfrac{16}{3}$ J.}
\opt{D}{If body starts from $x=-2$, it can go to $x=3$ even if its initial kinetic energy is zero.}
\end{qoptions}
\qn{4} Four identical rods which have thermally insulated lateral surfaces are joined at point $A$. Points $B,C,D$ \& $E$ are connected to large reservoirs. If heat flows into the junction from point $B$ at rate of $1$ W and from point $C$ at $3$ W inside, flows out from $D$ at $5$ W, which relation(s) is/are correct for temperature of these points ?
\qfig{p1_phy_s1_q04_1}
\begin{qoptions}
\opt{A}{$T_A<T_E$}
\opt{B}{$T_B=T_C$}
\opt{C}{$T_C>T_D$}
\opt{D}{$T_B=T_E$}
\end{qoptions}
\qn{5} In a tank filled with water, a small hole is made on the side. The emerging water jet falls to a distance of $0.8$ m from the base as shown ($y$ is distance of hole from above) :-
\qfig{p1_phy_s1_q05_1}
\begin{qoptions}
\opt{A}{As the level of water falls, the range decreases.}
\opt{B}{As the level of water falls, the time of flight for water jet remains the same.}
\opt{C}{$y$ can be $0.8$ m.}
\opt{D}{$y$ can be $0.2$ m.}
\end{qoptions}
\qn{6} In a standard YDSE setup, the fringe width on the screen is $1.5$ mm. When a thin glass film is pasted in front of the upper slit, the fringe pattern shifts up. But it is seen that at a point $P$ above central maxima where intensity was one fourth the intensity at central maxima, intensity remains the same. There were no maxima between central maxima and point $P$ before film was introduced. What can be the thickness of the film ? Take $\mu=1.5$, $\lambda=450$ nm :-
\begin{qoptions}
\opt{A}{$3\times10^{-7}$m}
\opt{B}{$5\times10^{-7}$ m}
\opt{C}{$9\times10^{-7}$ m}
\opt{D}{$1.4\times10^{-6}$ m}
\end{qoptions}
\sect{SEC--2}
\qn{1} A charged particle having non zero velocity is subjected to certain conditions given in List-I, List-II gives possible trajectories of the particle. Match the conditions in List-I with the results in List-II. Given all the fields are non-zero.\\
\textbf{List-I}
\begin{qoptions}
\opt{P}{In only uniform electric field}
\opt{Q}{In only uniform magnetic field}
\opt{R}{In uniform magnetic and uniform electric field}
\opt{S}{Subjected to a net force of constant magnitude}
\end{qoptions}
\noindent\textbf{List-II}
\begin{qoptions}
\opt{1}{The path of the charged particle may be a straight line}
\opt{2}{The path of the charged particle may be a parabola}
\opt{3}{The path of the charged particle may be circular}
\opt{4}{The path of the charged particle may be a helical with uniform or non-uniform pitch.}
\opt{5}{None of these}
\end{qoptions}
\begin{qoptions}
\opt{A}{P $\to$ 1,2; Q $\to$ 1,3,4; R $\to$ 1,4; S $\to$ 1,2,3,4}
\opt{B}{P $\to$ 1,2,3; Q $\to$ 1,3,4; R $\to$ 1,4; S $\to$ 1,2,3}
\opt{C}{P $\to$ 1,2; Q $\to$ 1,2,3,4; R $\to$ 1,3,4; S $\to$ 1,2,3,4}
\opt{D}{P $\to$ 1,2,3; Q $\to$ 1,3; R $\to$ 1,3,4; S $\to$ 1,2,3}
\end{qoptions}
\qn{2} In the list--I, there are certain situations depicted and list--II gives the description of speed and acceleration of particle. Match them. Neglect air resistance and assume uniform gravitational field wherever applicable. Consider time before it hits the ground.\\
\textbf{List-I}
\begin{qoptions}
\opt{P}{\qfigi{p1_phy_s2_q02_2}\\Particle moving in $1$ dimension}
\opt{Q}{$x=t^2+3t$}
\opt{R}{A projectile thrown on level ground at an angle to the horizontal}
\opt{S}{A ball thrown vertically down from a cliff.}
\end{qoptions}
\noindent\textbf{List-II}
\begin{qoptions}
\opt{1}{speed increases \& then decreases}
\opt{2}{speed decreases \& then increases}
\opt{3}{speed increases continuously.}
\opt{4}{speed decreases continuously}
\opt{5}{acceleration is initially opposite to velocity.}
\opt{6}{acceleration is initially in the same direction as the velocity}
\end{qoptions}
\begin{qoptions}
\opt{A}{P $\to$ 1,5; Q $\to$ 2,6; R $\to$ 4; S $\to$ 2,5}
\opt{B}{P $\to$ 2,5; Q $\to$ 3,6; R $\to$ 2; S $\to$ 3,6}
\opt{C}{P $\to$ 1,5; Q $\to$ 4,6; R $\to$ 1,6; S $\to$ 2,5}
\opt{D}{P $\to$ 2,6; Q $\to$ 2; R $\to$ 3,6; S $\to$ 3,6}
\end{qoptions}
\qn{3} An ideal gas is taken along the reversible processes as represented by the adjoining diagram.
\qfig{p1_phy_s2_q03_1}
\textbf{List-I}
\begin{qoptions}
\opt{P}{For process $B\to C$}
\opt{Q}{For process $A\to B$}
\opt{R}{For cycle $A\to B\to C\to A$}
\opt{S}{For process $C\to A$}
\end{qoptions}
\noindent\textbf{List-II}
\begin{qoptions}
\opt{1}{$\Delta Q>0$}
\opt{2}{$\Delta W>0$}
\opt{3}{$\Delta U>0$}
\opt{4}{$\Delta W=0$}
\opt{5}{$\Delta Q<0$}
\end{qoptions}
\begin{qoptions}
\opt{A}{P $\to$ 1,2,3; Q $\to$ 1,2,5; R $\to$ 1,2,3; S $\to$ 1,2}
\opt{B}{P $\to$ 1,2; Q $\to$ 1,3,4; R $\to$ 1,4,5; S $\to$ 5}
\opt{C}{P $\to$ 1,2,3; Q $\to$ 1,3,4; R $\to$ 1,2; S $\to$ 5}
\opt{D}{P $\to$ 1,3; Q $\to$ 1,3,4; R $\to$ 5; S $\to$ 4,5}
\end{qoptions}
\qn{4} Match the list--1 with nature of image in list--II (Consider all optical devices are kept in air)\\
\textbf{List-I}
\begin{qoptions}
\opt{P}{\qfigi{p1_phy_s2_q04_1}}
\opt{Q}{\qfigi{p1_phy_s2_q04_2}}
\opt{R}{\qfigi{p1_phy_s2_q04_3}}
\opt{S}{\qfigi{p1_phy_s2_q04_4}}
\end{qoptions}
\noindent\textbf{List-II}
\begin{qoptions}
\opt{1}{Erect image}
\opt{2}{Inverted image}
\opt{3}{magnified image}
\opt{4}{diminished image}
\opt{5}{real image}
\opt{6}{virtual image}
\end{qoptions}
\begin{qoptions}
\opt{A}{P $\to$ 2,3,5; Q $\to$ 1,4,6; R $\to$ 1,3,6; S $\to$ 2,3,5}
\opt{B}{P $\to$ 1,4; Q $\to$ 2,3,5; R $\to$ 1,4,6; S $\to$ 1,3,6}
\opt{C}{P $\to$ 1,4,6; Q $\to$ 1,3,6; R $\to$ 2,3,5; S $\to$ 1,4}
\opt{D}{P $\to$ 2,3,5; Q $\to$ 1,3,6; R $\to$ 1,3,6; S $\to$ 1,4,6}
\end{qoptions}
\sect{SEC--3}
\qn{1} The side of a cube is measured by vernier callipers ($10$ divisions of vernier scale coincide with $9$ divisions of main scale, where $1$ division of main scale is $1$ mm). The main scale reads $10$ mm and first division of vernier scale coincides with the main scale. Mass of the cube is $2.736$ g. Find the density (in g/cm$^3$) of the cube in appropriate significant figures.
\qn{2} Three uniformly charged infinite wires with linear charge density $\lambda$ are placed along $x$, $y$ and $z$ axis respectively. Find the flux of electric field through Gaussian surface given by $x^2+y^2+z^2=1$ ; $x>0$ ; $y>0$ ; $z>0$. If your answer is $\dfrac{n\lambda}{12\epsilon_0}$ fill the value of n.
\qn{3} When a resistance $R$ is connected in series with an element $A$, the electric current is found to be lagging behind the voltage by angle $30^\circ$. When the same resistance is connected in series with element $B$, current leads by $60^\circ$. When $R$, $A$, $B$ are connected in series, the current now leads voltage by $\theta$ which is equal to $\tan^{-1}\left(K/\sqrt3\right)$, then the value of $K$ is (assume same AC source is used in all cases)
\qn{4} An x-ray beam of monochromatic photon are incident on a metallic surface having a negligible work function (take $\phi=0$). It is seen that the wavelength of most energetic photoelectrons is equal to the wavelength of x-ray photons. Find the wavelength (in pm). Round off to nearest integer.
\qn{5} Consider a circuit with an alternating source and contains inductor and capacitor. Given reading of $A_1$ and $A_2$ as $3$ ampere and $5$ ampere respectively. Find the reading of $A$ in ampere.
\qfig{p1_phy_s3_q05_1}
\qn{6} In the circuit diagram shown, all the capacitors are in $\mu$F. The equivalent capacitance between points $A$ and $B$ is (in $\mu$ F).
\qfig{p1_phy_s3_q06_1}
\qn{7} A sample of $\beta$ active nuclei has a half life of $100$ min. Initially, there are $10^{13}$ $\beta$ active nuclei. Assume that all $\beta$ particles emitted leave the sample. What is the charge acquired (in $\mu$C) by the sample in $200$ min?
\qn{8} A cube is immersed in kerosene and water at rest without any external force as shown. To keep it completely submerged in water, we need to apply an extra force of $4$ N. If it was initially submerged to depth of $\dfrac a4$ in kerosene, what is the length `a' of the side of the cube (in cm) ?
\qfig{p1_phy_s3_q08_1}
\partsec{2}{CHEMISTRY}
\sect{SEC--1}
\qn{1} Which statement(s) is/are correct -
\begin{qoptions}
\opt{A}{If $0.1$M CH$_3$COOH aqueous solution ($K_a=1.8\times10^{-5}$) is diluted at $25^\circ$C, then [H$^+$] will increase.}
\opt{B}{On heating a sample of pure water [H$^+$] and [OH$^{-}$] increases but the sample still remains neutral}
\opt{C}{$10^{-4}$M HCl solution is less acidic than $0.1$M HCN ($K_{a\,\mathrm{HCN}}=10^{-5}$) solution}
\opt{D}{$10^{-8}$ M HCl aqueous solution will have pH exactly equal to $8$}
\end{qoptions}
\qfullbegin
\qn{2}
\qfig{p1_chem_s1_q02_1}
\qfullsplit
correct statements regarding the above reaction is/are :
\begin{qoptions}
\opt{A}{Free carbocation is involved in the reaction as an intermediate}
\opt{B}{Alkyl-oxygen bond cleavage takes place}
\opt{C}{Acyl-oxygen bond cleavage takes place}
\opt{D}{Alcohol formed in the reaction will give immediate turbidity with Lucas reagent}
\end{qoptions}
\qfullend
\qn{3} Which of the following reaction is/are feasible?
\begin{qoptions}
\opt{A}{$2[\mathrm{Fe(CN)_6}]^{3-}+2\mathrm{OH}^{-}+\mathrm{H_2O_2}\rightarrow2[\mathrm{Fe(CN)_6}]^{4-}+2\mathrm{H_2O}+\mathrm{O_2}$}
\opt{B}{$2\mathrm{K_4[Fe(CN)_6]}+\mathrm{H_2O_2}+\mathrm{H_2SO_4}\rightarrow2\mathrm{K_3[Fe(CN)_6]}+\mathrm{K_2SO_4}+2\mathrm{H_2O}$}
\opt{C}{$2\mathrm{KMnO_4}+3\mathrm{H_2O_2}\rightarrow2\mathrm{KOH}+2\mathrm{MnO_2}+2\mathrm{H_2O}+3\mathrm{O_2}$}
\opt{D}{$2\mathrm{KMnO_4}+5\mathrm{H_2O_2}+3\mathrm{H_2SO_4}\rightarrow2\mathrm{MnSO_4}+\mathrm{K_2SO_4}+5\mathrm{O_2}+8\mathrm{H_2O}$}
\end{qoptions}
\qn{4} Select the correct statement(s)?
\begin{qoptions}
\opt{A}{Addition of HgI$_2$ to aqueous solution of KI increases the freezing point since a complex K$_2$[HgI$_4$] is formed}
\opt{B}{Three aqueous solutions of ammonia labelled as A, B and C with concentrations $0.1$ M, $0.01$ M and $0.001$ M respectively, then value of Van't Hoff factor for these solution will be in order :- $i_A<i_B<i_C$ (consider degree of dilution)}
\opt{C}{Vapour density of PCl$_5$(g) dissociating into PCl$_3$ (g) and Cl$_2$(g) is $100$. \% dissociation of PCl$_5$ is $4.25\%$ [Given : atomic mass of P $=31$, Cl $=35.5$]}
\opt{D}{If pH of $0.1$ M monobasic acid is found to be $2$, its osmotic pressure at a given temperature $T$ K is $0.11$ RT}
\end{qoptions}
\qn{5} Compounds having heat of hydrogenation lesser than that of cyclohexa-1, 4-diene i.e. \qfigi{p1_chem_s1_q05_1}

[Hint resonance energy of benzene is $36$ Kcal]
\begin{qoptions}
\opt{A}{\qfigi{p1_chem_s1_q05_A}}
\opt{B}{\qfigi{p1_chem_s1_q05_B}}
\opt{C}{\qfigi{p1_chem_s1_q05_C}}
\opt{D}{\qfigi{p1_chem_s1_q05_D}}
\end{qoptions}
\qn{6} Which of the following compound/s does not liberate acidic gas during its hydrolysis.
\begin{qoptions}
\opt{A}{Ca$_3$P$_2$}
\opt{B}{AlN}
\opt{C}{Al$_2$S$_3$}
\opt{D}{CaH$_2$}
\end{qoptions}
\sect{SEC--2}
\qn{1} Match the list.\\
\textbf{List-I} (Electrolysis using inert electrode)
\begin{qoptions}
\opt{P}{Aqueous solution of sodium acetate}
\opt{Q}{Aqueous solution of sodium formate}
\opt{R}{Concentrated aqueous solution of NaCl}
\opt{S}{Dilute aqueous solution of HCl}
\end{qoptions}
\noindent\textbf{List-II} (Product at anode or cathode.)
\begin{qoptions}
\opt{1}{H$_2$}
\opt{2}{CO$_2$}
\opt{3}{Cl$_2$}
\opt{4}{C$_2$H$_6$}
\opt{5}{O$_2$}
\end{qoptions}
The correct option is -
\begin{qoptions}
\opt{A}{P $\to$ 1,2,4; Q $\to$ 1,2; R $\to$ 1, 3; S $\to$ 1, 5}
\opt{B}{P $\to$ 1,5; Q $\to$ 2,4; R $\to$ 1; S $\to$ 3, 4}
\opt{C}{P $\to$ 1,4; Q $\to$ 2,4; R $\to$ 1, 5; S $\to$ 3, 5}
\opt{D}{P $\to$ 2,4; Q $\to$ 1,4; R $\to$ 2, 5; S $\to$ 1,5}
\end{qoptions}
\qn{2} Match the list\\
\textbf{List-I}
\begin{qoptions}
\opt{P}{\qfigi{p1_chem_s2_q02_1}}
\opt{Q}{\qfigi{p1_chem_s2_q02_2}}
\opt{R}{\qfigi{p1_chem_s2_q02_3}}
\opt{S}{\qfigi{p1_chem_s2_q02_4}}
\end{qoptions}
\noindent\textbf{List-II}
\begin{qoptions}
\opt{1}{Ionic organometalic product is formed}
\opt{2}{Product can show electrophillic addition reaction with HBr (1 eq.)}
\opt{3}{Reaction involve free radical as an intermediate}
\opt{4}{Number of carbon in parent chain of product is $4$}
\end{qoptions}
Select correct code for your answer :
\begin{qoptions}
\opt{A}{P $\to$ 3, 4; Q $\to$ 2; R $\to$ 2, 4; S $\to$ 1, 4}
\opt{B}{P $\to$ 3, 4; Q $\to$ 2; R $\to$ 2, 4; S $\to$ 1, 3}
\opt{C}{P $\to$ 3, 4; Q $\to$ 1, 2; R $\to$ 2, 4; S $\to$ 1, 4}
\opt{D}{P $\to$ 2, 3; Q $\to$ 2; R $\to$ 2, 4; S $\to$ 1, 4}
\end{qoptions}
\qn{3} Match the list:\\
\textbf{List-I} (complex)
\begin{qoptions}
\opt{P}{[CoF$_6$]$^{3-}$}
\opt{Q}{[Cr(C$_2$O$_4$)$_3$]$^{3-}$}
\opt{R}{[AuCl$_4$]$^{-}$}
\opt{S}{[Fe(CN)$_6$]$^{3-}$}
\end{qoptions}
\noindent\textbf{List-II} (Properties)
\begin{qoptions}
\opt{1}{Diamagnetic and low spin complex}
\opt{2}{Paramagnetic and outer orbital complex}
\opt{3}{Paramagnetic and inner orbital complex}
\opt{4}{Paramagnetic and low spin complex}
\end{qoptions}
The correct option is
\begin{qoptions}
\opt{A}{P $\to$ 1; Q $\to$ 2; R $\to$ 3; S $\to$ 4}
\opt{B}{P $\to$ 2; Q $\to$ 3; R $\to$ 1; S $\to$ 4}
\opt{C}{P $\to$ 2; Q $\to$ 1; R $\to$ 3; S $\to$ 4}
\opt{D}{P $\to$ 1; Q $\to$ 3; R $\to$ 2; S $\to$ 4}
\end{qoptions}
\qn{4} Match the list:\\
\textbf{List-I} (Complex Compound)
\begin{qoptions}
\opt{P}{[Ni(H$_2$O)$_6$]Cl$_2$}
\opt{Q}{[Co(CN)$_2$(NH$_3$)$_4$]OC$_2$H$_5$}
\opt{R}{[IrCl$_6$]$^{3-}$}
\opt{S}{[PtCl$_2$(NH$_3$)$_4$]Br$_2$}
\end{qoptions}
\noindent\textbf{List-II} (Properties)
\begin{qoptions}
\opt{1}{d$^2$sp$^3$ hybridisation}
\opt{2}{Ionisation isomerism}
\opt{3}{$\mu=2.83$ B.M.}
\opt{4}{$\Delta_0<$ P}
\opt{5}{sp$^3$d$^2$ hybridisation}
\end{qoptions}
The correct option is
\begin{qoptions}
\opt{A}{P $\to$ 3,4,5; Q $\to$ 1,2; R $\to$ 1; S $\to$ 1,2}
\opt{B}{P $\to$ 3,4; Q $\to$ 1; R $\to$ 2,3; S $\to$ 1,2}
\opt{C}{P $\to$ 1,2,3; Q $\to$ 1,4; R $\to$ 4,5; S $\to$ 2,5}
\opt{D}{P $\to$ 2,3; Q $\to$ 2; R $\to$ 3,4; S $\to$ 1}
\end{qoptions}
\sect{SEC--3}
\qn{1} For the first order reaction A(g) $\to$ 2B(g) $+$ C(g) the half life for decomposition of A is $3$ minute at $300$K. Calculate the time (in minute) in which partial pressure of A(g) will drop from $2$ bar to $0.5$ bar at $400$K. Given activation energy of reaction is $840$ R. (Take $\ln2=0.7$)
\qn{2} $2$ moles of methane gas are compressed from ($1$ bar, $2$L) to $2$ bar isothermally against constant external pressure. Calculate work done in Joules.
\qn{3} For $10$ minutes each, at $0^\circ$C, from two identical holes nitrogen and an unknown gas are leaked into a common vessel of $4$ litre capacity. The resulting pressure in vessel is $2.8$ atm and the mixture contains $0.4$ mole of nitrogen. What is the molar mass of unknown gas? (Take R $=0.0821$ L--atm - mol$^{-1}$ K$^{-1}$)
\qn{4}
\qfig{p1_chem_s3_q04_1}
Find out total number of final products?
\qfullbegin
\qn{5} How many of the following compounds given below have faster aromatic nucleophilic substitution
\qfig{p1_chem_s3_q05_1}
\qfullend
\qn{6} Find out total number of moles of Grignard reagent RMgX consumed when treated it with $1$ mole of `X'
\qfig{p1_chem_s3_q06_1}
\qn{7} Find the number of electrons having $|m|\le1$, $l=2$, $m_s=-\dfrac12$ in Pd ($Z=46$).
\qn{8} During extraction of iron, CO is the chief reducing agent. If $72$ Kg of carbon is available and produced CO is completely used to reduce Haematite (present in excess amount) then calculate how much Kg of Iron is obtained after reduction. (Given : atomic weight of C--12 g/mol, O--16 g/mol, Fe--55.85 g/mol)
\partsec{3}{MATHEMATICS}
\sect{SEC--1}
\qn{1} If $e_1$ is eccentricity of $9x^2+4y^2=36$ and $e_2$ is eccentricity of $9x^2-4y^2=36$ then
\begin{qoptions}
\opt{A}{$e_1^2+e_2^2<3$}
\opt{B}{$e_1^2+e_2^2>3$}
\opt{C}{$e_1^2+e_2^2<4$}
\opt{D}{$e_1^2+e_2^2>4$}
\end{qoptions}
\qn{2} Let $f:[0,1]\to[0,\infty)$ be a differentiable function so that $f'(x)$ is decreasing function ; $f(0)=0$, $f'(1)>0$, $f'(0)>0$ and $I=\displaystyle\int_0^1\dfrac{dx}{1+f^2(x)}\,dx$ then which of the following is \textbf{INCORRECT} ?
\begin{qoptions}
\opt{A}{$I\le\dfrac{f(1)}{f'(1)}$}
\opt{B}{$I<\dfrac{f(1)}{f'(1)}$}
\opt{C}{$I\ge\dfrac{f(1)}{f'(1)}$}
\opt{D}{$I>\dfrac{f(1)}{f'(1)}$}
\end{qoptions}
\qn{3} The coordinates of four angular points of a tetrahedron $OABC$ are $(0,0,0)$, $(0,0,2)$, $(0,4,0)$ and $(6,0,0)$ respectively. A point $P$ inside the tetrahedron is at the same distance $r$ from the four plane faces of tetrahedron. Which of following \textbf{cannot} be the value of $r$ ?
\begin{qoptions}
\opt{A}{$\dfrac23$}
\opt{B}{$\dfrac25$}
\opt{C}{$2$}
\opt{D}{$3$}
\end{qoptions}
\qn{4} Let $C_k={}^nC_k$ for $0\le k\le n$ and $A_k=\begin{bmatrix}C_{k-1}&0\\0&C_k\end{bmatrix}$ for $k\ge1$ and $A_1A_2+A_2A_3+A_3A_4+\ldots+A_{n-1}A_n=\begin{bmatrix}a&0\\0&b\end{bmatrix}$ then
\begin{qoptions}
\opt{A}{$a=b$}
\opt{B}{$a+b=2\left({}^{2n}C_{n-1}-n\right)$}
\opt{C}{$a=2b$}
\opt{D}{$a+b=0$}
\end{qoptions}
\qn{5} Probability that the length of a randomly chosen chord of a circle lies between $\dfrac12$ and $\dfrac34$ of its diameter is $\dfrac pq$ where $p$ and $q$ are coprime number, then
\begin{qoptions}
\opt{A}{$p+q=21$}
\opt{B}{$q-p=11$}
\opt{C}{$p+q=20$}
\opt{D}{$q-p=10$}
\end{qoptions}
\qn{6} For two complex numbers $z_1$ and $z_2$ ; $\left(az_1+b\bar z_1\right)\left(cz_2+d\bar z_2\right)=\left(cz_1+d\bar z_1\right)\left(az_2+b\bar z_2\right)$ if
\begin{qoptions}
\opt{A}{$\dfrac ab=\dfrac cd$}
\opt{B}{$\dfrac ad=\dfrac bc$}
\opt{C}{$|z_1|=|z_2|$}
\opt{D}{$\arg(z_1)=\arg(z_2)$}
\end{qoptions}
\sect{SEC--2}
\qn{1} Let $\phi(x)=3f\left(\dfrac{x^2}{3}\right)+f\left(3-x^2\right)\ \forall\ x\in(-3,4)$ where $f''(x)>0\ \forall\ x\in(-3,4)$, $g(x)=\left(a^2-3a+2\right)\left(\cos^2\dfrac x4-\sin^2\dfrac x4\right)+(a-1)x+\sin1$\\
Match List-I with List-II and select the correct answer using the code given below the list.\\
\textbf{List-I}
\begin{qoptions}
\opt{P}{$\phi(x)$ is increasing in}
\opt{Q}{$\phi(x)$ is decreasing in}
\opt{R}{The set of values of $a$ for which $g(x)$ does not possess critical point is}
\opt{S}{The set of values of $a$ for which $g(x)$ possess critical point is}
\end{qoptions}
\noindent\textbf{List-II}
\begin{qoptions}
\opt{1}{$(0,1)\cup(1,4)$}
\opt{2}{$\left[-\dfrac32,0\right]\ ;\ \left[\dfrac32,4\right)$}
\opt{3}{$\left(-3,-\dfrac32\right]\ ;\ \left[0,\dfrac32\right]$}
\opt{4}{$(-\infty,0]\cup\{1\}\cup[4,\infty)$}
\end{qoptions}
\begin{qoptions}
\opt{A}{P $\to$ 2; Q $\to$ 3; R $\to$ 1; S $\to$ 4}
\opt{B}{P $\to$ 2; Q $\to$ 3; R $\to$ 4; S $\to$ 1}
\opt{C}{P $\to$ 3; Q $\to$ 2; R $\to$ 1; S $\to$ 4}
\opt{D}{P $\to$ 3; Q $\to$ 2; R $\to$ 4; S $\to$ 1}
\end{qoptions}
\qn{2} $S\equiv x^2+y^2-2x-4y+1=0$ and $S'\equiv x^2+y^2-4x-2y+1=0$ are equations of two circles.\\
Match List-I with List-II and select the correct answer using the code given below the list.\\
\textbf{List-I}
\begin{qoptions}
\opt{P}{Sum of number of tangents from $(1,3)$ to $S$ and number of tangents from $(1,3)$ to $S'$ is}
\opt{Q}{If $r$ is radius of circle passes through point of intersection of $S$ and $S'$ and whose normal is the straight line $2x+y-6=0$, then the value of $\dfrac{r^2}{4}$ is}
\opt{R}{If $m$ is slope of common tangent of $S=0$ and $S'=0$ then $|m|$ is}
\opt{S}{If $\ell$ is length of common tangent of $S$ and $S'$ then the value of $2\ell^2$ is}
\end{qoptions}
\noindent\textbf{List-II}
\begin{qoptions}
\opt{1}{$1$}
\opt{2}{$2$}
\opt{3}{$3$}
\opt{4}{$4$}
\end{qoptions}
\begin{qoptions}
\opt{A}{P $\to$ 2; Q $\to$ 1; R $\to$ 2; S $\to$ 3}
\opt{B}{P $\to$ 2; Q $\to$ 2; R $\to$ 1; S $\to$ 4}
\opt{C}{P $\to$ 2; Q $\to$ 1; R $\to$ 4; S $\to$ 1}
\opt{D}{P $\to$ 1; Q $\to$ 1; R $\to$ 4; S $\to$ 2}
\end{qoptions}
\qn{3} (1) The sum of all the values of $r$ satisfying ${}^{39}C_{3r-1}-{}^{39}C_{r^2}={}^{39}C_{r^2-1}-{}^{39}C_{3r}$ is $\alpha_1$.\\
(2) If ${}^{2n+3}C_{2n}-{}^{2n+2}C_{2n-1}=15(2n+1)$ then the value of $n$ is $\alpha_2$\\
(3) If ${}^{56}P_{r+6}:{}^{54}P_{r+3}=30800:1$ then the value of $r$ is $\alpha_3$\\
(4) ${}^{n+2}C_8:{}^{n-2}P_4=57:16$ then the value of $n$ is $\alpha_4$\\
Match List-I with List-II and select the correct answer using the code given below the list.\\
\textbf{List-I}
\begin{qoptions}
\opt{P}{The value of $\alpha_1$ is}
\opt{Q}{The value of $\alpha_2$ is}
\opt{R}{The value of $\alpha_3$ is}
\opt{S}{The value of $\alpha_4$ is}
\end{qoptions}
\noindent\textbf{List-II}
\begin{qoptions}
\opt{1}{$41$}
\opt{2}{$8$}
\opt{3}{$14$}
\opt{4}{$19$}
\end{qoptions}
\begin{qoptions}
\opt{A}{P $\to$ 2; Q $\to$ 3; R $\to$ 1; S $\to$ 4}
\opt{B}{P $\to$ 2; Q $\to$ 3; R $\to$ 4; S $\to$ 1}
\opt{C}{P $\to$ 2; Q $\to$ 1; R $\to$ 3; S $\to$ 4}
\opt{D}{P $\to$ 3; Q $\to$ 2; R $\to$ 1; S $\to$ 4}
\end{qoptions}
\qn{4} (i) Let harmonic mean, arithmetic mean and geometric mean of two positive numbers $a$ and $b$ respectively be $4$, $A$ and $G$. If $2A+G^2=27$ then $a^2+b^2=\alpha_1$\\
(ii) If the value of $x+y+z=15$ where $a,x,y,z,b$ are in arithmetic progression while the value of $\dfrac1x+\dfrac1y+\dfrac1z$ is $\dfrac58$ where $a,x,y,z,b$ are in harmonic progression, then the value of $a^2+b^2$ is $\alpha_2$\\
(iii) The value of $\displaystyle\sum_{r=1}^{16}\left(\dfrac{1^3+2^3+\ldots..+r^3}{1+3+\ldots..+(2r-1)}\right)$ is $\alpha_3$\\
(iv) If the geometric mean and harmonic mean of two positive numbers $x_1$ and $x_2$ are $18$ and $\dfrac{216}{13}$ respectively then the value of $|x_1-x_2|$ is $\alpha_4$.\\
Match List-I with List-II and select the correct answer using the code given below the list.\\
\textbf{List-I}
\begin{qoptions}
\opt{P}{The value of $\alpha_1$ is}
\opt{Q}{The value of $\alpha_2$ is}
\opt{R}{The value of $\alpha_3$ is}
\opt{S}{The value of $\alpha_4$ is}
\end{qoptions}
\noindent\textbf{List-II}
\begin{qoptions}
\opt{1}{$52$}
\opt{2}{$45$}
\opt{3}{$446$}
\opt{4}{$15$}
\end{qoptions}
\begin{qoptions}
\opt{A}{P $\to$ 2; Q $\to$ 1; R $\to$ 2; S $\to$ 3}
\opt{B}{P $\to$ 2; Q $\to$ 2; R $\to$ 1; S $\to$ 4}
\opt{C}{P $\to$ 2; Q $\to$ 1; R $\to$ 3; S $\to$ 4}
\opt{D}{P $\to$ 1; Q $\to$ 1; R $\to$ 4; S $\to$ 2}
\end{qoptions}
\sect{SEC--3}
\qn{1} In a triangle $ABC$ (with usual notations), $a=5$, $b=4$ and $A-B=\cos^{-1}\left(\dfrac{31}{32}\right)$ then the value of $\dfrac c5$ is
\qn{2} In the expansion of $(1-x)^{30}(1+x+x^2)^{31}$ the coefficient of $x^{47}$ is $(-{}^{30}C_r)$ then the value of $\left|\dfrac r2\right|$ is
\qn{3} In the equation $x^5-5x^4+9x^3-9x^2+5x-1=0$ if $p$ is the sum of real roots and $q$ is the product of real roots then the value of $\dfrac qp$ is
\qn{4} Let $f:(0,\infty)\to(0,\infty)$ be a function differentiable at $3$ and satisfying $f(3)=3f'(3)$. If the value of $\displaystyle\lim_{x\to\infty}\left(\dfrac{f\left(3+\frac3x\right)}{f(3)}\right)^x$ is $L$ then $\dfrac{[L]}{5}$ is (where $[.]$ is greatest integer function)
\qn{5} The value of $\displaystyle\int_{\pi/4}^{3\pi/4}\dfrac{\cos x}{1-\cos x}\,dx$ is equal to (use $\pi=\dfrac{22}{7}$)
\qn{6} Let
\[f(x)=\begin{vmatrix}\sec x&\cos x&\sec^2x+\cot x\operatorname{cosec}x+\cos x\\\cos^2x&\cos^2x&\operatorname{cosec}^2x+\cos^4x\\1&\cos^2x&\cos^2x\end{vmatrix}\]
then the value $\left|\displaystyle\int_0^{\pi/2}f(x)\,dx\right|$ is (use $\pi=\dfrac{22}{7}$)
\qn{7} If $y=2x-3$ is a tangent to the parabola $y^2=4a\left(x-\dfrac13\right)$ then the value of $|a|$ is
\qn{8} $OA$, $OB$, $OC$ are the sides of a rectangular parallelopiped whose diagonals are $OO'$, $AA'$, $BB'$ and $CC'$. $D$ is the centre of rectangle $AC'O'B$ and $D'$ is centre of rectangle $O'A'CB'$. If sides $OA$, $OB$, $OC$ are in ratio $1:2:3$ and $\angle DOD'=\alpha$ then $\dfrac{697\cos^2\alpha}{128}$ is\\
(Centre of rectangle is mid point of diagonal)
\papersec{2}
\partsec{1}{PHYSICS}
\sect{SEC--1}
\qn{1} Pick the correct option/options.
\begin{qoptions}
\opt{A}{If the sum of all the forces acting on a particle is zero, then it will necessarily be in equilibrium.}
\opt{B}{When a sphere rolls on a horizontal fixed surface, there is no point on the sphere which has vertical velocity.}
\opt{C}{When tall buildings are constructed on earth, the duration of day and night increases.}
\opt{D}{A body is rotating non uniformly about a vertical axis fixed in an inertial reference frame. The resultant force on the particle of the body not on the axis is horizontal and skew with the axis.}
\end{qoptions}
\qn{2} A charged ball of charge $+q_1$ is revolving around another charge $+q$ as shown in a conical pendulum. The motion is in a horizontal plane :-
\qfig{p2_phy_s1_q02_1}
\begin{qoptions}
\opt{A}{Tension in the string is greater than the weight of the ball.}
\opt{B}{The tension in the string is greater than the electrostatic repulsive force.}
\opt{C}{If the charge is removed, the speed of ball has to be increased to maintain the angle.}
\opt{D}{If the charge is removed, the speed of ball has to be decreased to maintain the angle.}
\end{qoptions}
\qn{3} Mark the incorrect statement(s) for a standing wave established in a string :
\begin{qoptions}
\opt{A}{Energy of an element of medium remains constant.}
\opt{B}{Phase difference between any two particles can be $\dfrac\pi2$}
\opt{C}{Energy can get transferred through antinodes.}
\opt{D}{Energy can get transferred through nodes.}
\end{qoptions}
\qn{4} A wall is moving with velocity $u$ and a source of sound (frequency $f$) moves with velocity $\dfrac u2$ in the same direction as shown. Assuming that the sound travels with a speed of $10u$, pick the right option or options.
\qfig{p2_phy_s1_q04_1}
\begin{qoptions}
\opt{A}{wavelength of incident sound is $\dfrac{19u}{2f}$.}
\opt{B}{frequency of the incident sound is $\dfrac{18}{19}f$.}
\opt{C}{wavelength of reflected sound is $\dfrac uf$.}
\opt{D}{frequency of the reflected sound is $\dfrac{12f}{19}$.}
\end{qoptions}
\qn{5} A gas is undergoing an adiabatic process. At a certain stage, the values of volume and temperature are $V_0$ and $T_0$ respectively and the modulus of slope of $V$--$T$ curve is $m$. Choose the correct options-
\begin{qoptions}
\opt{A}{The value of $C_V=\dfrac{RT_0m}{V_0}$}
\opt{B}{The value of $C_P=\left(1+\dfrac{T_0m}{V_0}\right)R$}
\opt{C}{$m=\dfrac{V_0}{T_0(\gamma-1)}$}
\opt{D}{$m=\dfrac{-V_0}{T_0(\gamma-1)}$}
\end{qoptions}
\qn{6} A long straight cylindrical shell has inner radius $R_i$ and outer radius $R_0$. It carries current $i$, uniformly distributed over its cross-section. A solid wire of radius $r$ is to be placed parallel to the cylinder axis, in the hollow region such that ($r<R_i$) :-
\begin{qoptions}
\opt{A}{The magnetic field is zero everywhere in hollow region. We conclude that the wire is on the cylinder axis and carries current $i$ in the same direction as current in the shell.}
\opt{B}{The magnetic field is zero everywhere outside the shell ($r>R_0$). We conclude that the wire is on the cylinder axis and carries $i$ in the direction opposite to the direction of current in the shell.}
\opt{C}{The magnetic field is zero everywhere in the hollow region. We conclude that the wire may be anywhere in the hollow region but must be carrying current $i$ in the same direction as the current in the shell.}
\opt{D}{The magnetic field is zero everywhere in the hollow region. We conclude that wire does not carry any current.}
\end{qoptions}
\qn{7} The diagram shows a modified meter bridge, which is used for measuring two unknown resistance at the same time. When only the first galvanometer $P$ is used, for obtaining the balance point, it is found at point $C$. Now the first galvanometer $P$ is removed and the second galvanometer $Q$ is used, which gives balance point at $D$. Using the details given in the diagram, find out the value of $R_1$ and $R_2$.
\qfig{p2_phy_s1_q07_1}
\begin{qoptions}
\opt{A}{$R_1=5R/3$}
\opt{B}{$R_2=4R/3$}
\opt{C}{$R_1=4R/3$}
\opt{D}{$R_2=5R/3$}
\end{qoptions}
\qn{8} A mirror is hanging by a light string of length $\ell$ from ceiling. A light pulse of monochromatic light rays is incident normally on the mirror and gets reflected. Which of the following statements is/are correct ?
\begin{qoptions}
\opt{A}{Mirror will start to oscillate.}
\opt{B}{Wavelength of reflected rays will be greater than that of the incident rays.}
\opt{C}{Wavelength of reflected rays may be less than that of incident rays.}
\opt{D}{Assuming the intensity of incident light to be constant, the maximum angle to which the mirror deflects depends on the wavelength of the light pulse.}
\end{qoptions}
\sect{SEC--2}
\paraq{1}{2}
For given circuit, give answer of following questions
\qfig{p2_phy_s2_para1_2_1}
\qn{1} If device--1 is battery of emf $10$V and $2\Omega$ internal resistance (+ive end at $A$) and device-2 is battery of emf $x$ volt and internal resistance $1\Omega$ (+ive end at $C$) and reading of galvanometer is zero then value of $x$ is :-
\begin{qoptions}
\opt{A}{$6$ volt}
\opt{B}{$4$ volt}
\opt{C}{$10$ volt}
\opt{D}{$2$ volt}
\end{qoptions}
\qn{2} If device-1 and device-2 are capacitors of capacitance $C_1$ and $C_2$ and reading of galvanometer remain zero then $\dfrac{C_1}{C_2}$ is (no charge pass through Galvanometer):-
\begin{qoptions}
\opt{A}{$\dfrac13$}
\opt{B}{$\dfrac12$}
\opt{C}{$1$}
\opt{D}{$\dfrac23$}
\end{qoptions}
\paraq{3}{4}
A charged particle having mass `$m$' and charge `$+q$' is released from rest at origin in uniform electric field $\vec E=E_0\hat j$ and uniform magnetic field $\vec B=B_0\hat k$.
\qn{3} Maximum value of `$y$' coordinate of particle in the path.
\begin{qoptions}
\opt{A}{$\dfrac{mE_0}{qB_0^2}$}
\opt{B}{$\dfrac{2mE_0}{qB_0^2}$}
\opt{C}{$\dfrac{3mE_0}{qB_0^2}$}
\opt{D}{$\dfrac{4mE_0}{qB_0^2}$}
\end{qoptions}
\qn{4} Mark the correct option(s) :
\begin{qoptions}
\opt{A}{When $y$-coordinate is maximum, the speed of the particle is maximum.}
\opt{B}{When $y$-coordinate is maximum, the velocity of particle is in $x$-direction.}
\opt{C}{Radius of curvature of path when $y$-coordinate is maximum is $\dfrac{4mE_0}{qB_0^2}$.}
\opt{D}{Radius of curvature of path when $y$-coordinate is maximum is $\dfrac{2mE_0}{qB_0^2}$.}
\end{qoptions}
\sect{SEC--3}
\qn{1} A uniform solid hemisphere of mass $m$ and radius $R$ is attached to the roof with a chord of torsional constant $C$ and performing torsional SHM. Then the time period (in seconds) of SHM is (Take $m=15$ kg, $R=\dfrac2\pi$ m, $C=6$ Nm/rad).
\qfig{p2_phy_s3_q01_1}
\qn{2} Young's double slit experiment is conducted in a liquid of refractive index $\mu_1$ as shown in figure. A thin transparent slab of refractive index $\mu_2$ is placed in front of the slit $S_2$. If $\mu_2=\dfrac32$, $\mu_1=\dfrac43$, $t=8$ $\mu$m, the magnitude of optical path difference (wrt liquid) at `O' in $\mu$m is :
\qfig{p2_phy_s3_q02_1}
\qn{3} We would like to increase the length of a $15$ cm long copper rod of cross-section $4$ mm$^2$ by $1$mm. The energy absorbed by the rod if it is heated is $E_1$. The energy absorbed by the rod if it is stretched slowly is $E_2$. The ratio of $E_1$ and $E_2$ is $\alpha$. The value of $\dfrac\alpha{100}$ is. [Various parameters of Copper are: Density $=9\times10^3$ kg/m$^3$, Thermal coefficient of linear expansion $=16\times10^{-6}$ K$^{-1}$, Young's modulus $=135\times10^9$ Pa, Specific heat $=400$ J/kg-K]
\qn{4} A pendulum bob hinged at $O$ is given velocity $u=\sqrt{4g\ell}$ from bottommost point as shown. If $\ell=54$ cm, find maximum height (in m) reached by bob from point of projection.
\qfig{p2_phy_s3_q04_1}
\qn{5} A concave mirror of focal length $20$ cm and a convex lens of focal length $10$ cm are kept with their optic axis parallel but separated by $0.5$ mm as shown in figure. The distance between lens and mirror is $10$ cm. An object of height $3$ mm is placed on the optic axis of lens at distance $15$ cm from the lens. Find height of image formed by mirror after refraction from lens (in mm).
\qfig{p2_phy_s3_q05_1}
\qn{6} A uniform disc of radius $R$ having charge $Q$ distributed uniformly all over its surface is placed on a smooth horizontal surface. A magnetic field $B=Kxt^2$, where $K=$ constant, $x$ is the distance (in metre) from the centre of the disc and $t$ is the time (in second) is switched on perpendicular to the plane of the disc. The torque (in N-m) acting on the disc after $15$ sec. (Take $2KQ=1$ S.I. unit and $R=1$ metre).
\qn{7} A small ring $R$ of mass $0.2$ kg is threaded onto a fixed rod which is inclined at an angle of $30^\circ$ to the horizontal. The ring is attached to a block $B$ of mass $0.5$ kg by means of a light string which passes over a smooth pulley $P$. The part $RP$ of the string is perpendicular to the rod, and the part $PB$ of the string is vertical (see diagram). The system is in equilibrium. Assuming the equilibrium is limiting, if the coefficient of friction ($\mu$) between $R$ and the rod is $\mu=\dfrac{1}{\alpha-\sqrt\beta}$ where $\alpha$ and $\beta$ are the smallest possible integer. Fill the value of ($\alpha+\beta$).
\qfig{p2_phy_s3_q07_1}
\qn{8} Two uniform rods of same length but different density are arranged as shown in the diagram. Centre of mass of this arrangement lies on a line $ax+by=\ell$. ($\ell$ is the length of rod) find $a+b=$ ?
\qfig{p2_phy_s3_q08_1}
\partsec{2}{CHEMISTRY}
\sect{SEC--1}
\qn{1} For $0.1$ M aq. solution of Na$_2$CO$_3$ at $25^\circ$C, select the correct options (Use : H$_2$CO$_3$ : $K_{a_1}=10^{-5}$, $K_{a_2}=10^{-9}$ ) and $K_w$ (H$_2$O) $=10^{-14}$
\begin{qoptions}
\opt{A}{pH of the solution is $11$.}
\opt{B}{Equilibrium concentration of H$_2$CO$_3$ is $10^{-9}$ M.}
\opt{C}{Equilibrium concentration of HCO$_3^-$ is $10^{-3}$ M}
\opt{D}{Equilibrium concentration of CO$_3^{2-}$ is $0.01$ M}
\end{qoptions}
\qn{2} Select the incorrect options :
\begin{qoptions}
\opt{A}{Gold sol is a positive charged sol.}
\opt{B}{Lyophobic sols can act as protective colloid for Lyophilic sols.}
\opt{C}{When FeCl$_3$ is hydrolysed with hot water, then Ferric hydroxide sol (positively charged) will be obtained.}
\opt{D}{Blood is negatively charged sol.}
\end{qoptions}
\qfullbegin
\qn{3}
\qfig{p2_chem_s1_q03_1}
\qfullsplit
Select the correct statements(s) for above reaction sequence ?
\begin{qoptions}
\opt{A}{F and G are identical.}
\opt{B}{E and H are optically inactive due to external compensation.}
\opt{C}{Compound (A) and compound (B) can be differentiated by Tollen's reagent}
\opt{D}{F and G are enantiomers}
\end{qoptions}
\qfullend
\qn{4} Among the following pair of compounds, identify the pair(s) which can be differentiated by Fehling as well as Tollen's reagent ?
\begin{qoptions}
\opt{A}{\qfigi{p2_chem_s1_q04_A}}
\opt{B}{\qfigi{p2_chem_s1_q04_B}}
\opt{C}{\qfigi{p2_chem_s1_q04_C}}
\opt{D}{\qfigi{p2_chem_s1_q04_D}}
\end{qoptions}
\qfullbegin
\qn{5} Compound (A) has molecular formula C$_6$H$_7$N, undergoes following reaction sequence.
\qfig{p2_chem_s1_q05_1}
\qfullsplit
all are major organic product in above reaction sequence. Select the correct statement(s) for the reaction sequence.
\begin{qoptions}
\opt{A}{Order of basic strength is A $>$ D $>$ E .}
\opt{B}{Compound (A) on reaction with CHCl$_3$, Alc. KOH results formation of \qfigi{p2_chem_s1_q05_2}}
\opt{C}{Reactivity order towards aromatic electrophilic substitution is A $>$ B.}
\opt{D}{Compound (E) on reaction with SnCl$_2$, HCl followed by H$_3$O$^+$ gives \qfigi{p2_chem_s1_q05_3}}
\end{qoptions}
\qfullend
\qn{6} Which of the following is/are true for the complex compound with IUPAC name as : Dichloridobis(ethane-1, 2-diamine)cobalt(III) sulphate
\begin{qoptions}
\opt{A}{Its formula is $\left[\mathrm{Co}\left(\mathrm{NH_2CH_2CH_2NH_2}\right)_2\mathrm{Cl_2}\right]\mathrm{SO_4}$}
\opt{B}{The co-ordination number of central metal ion is $6$}
\opt{C}{The effective atomic number of central metal is $32$}
\opt{D}{The compound can be categorised as a heteroleptic complex}
\end{qoptions}
\qn{7} The compound(s) that give more than one gaseous product on strong heating with Bunsen burner is/are
\begin{qoptions}
\opt{A}{Baking soda}
\opt{B}{Washing soda}
\opt{C}{Chile salt petre}
\opt{D}{Calcium oxalate}
\end{qoptions}
\qn{8} The reaction Fe $+$ dilH$_2$SO$_4$ $\to$ ; gives the product(s)
\begin{qoptions}
\opt{A}{FeSO$_4$}
\opt{B}{Fe$_2$(SO$_4$)$_3$}
\opt{C}{H$_2$}
\opt{D}{SO$_2$}
\end{qoptions}
\sect{SEC--2}
\paraq{1}{2}
For reversible Galvanic cell, thermodynamic parameters are given as :
\begin{gather*}
(\Delta G)_{T,P}=\Delta H-T\Delta S\\
(\Delta G)_{T,P}=-nFE_{\mathrm{cell}}\\
\Delta S=nF\left[\dfrac{dE_{\mathrm{cell}}}{dT}\right]_P\\
W_{\mathrm{usefull}}=-(\Delta G)_{T,P}
\end{gather*}
Answer the following questions :
\qn{1} For given cell Pt $|$ H$_2$ $|$ HCl $|$ AgCl $|$ Ag, variation of e.m.f. of the cell with temperature is given as $E_{\mathrm{cell}}$ (in volt) $=0.4-2\times10^{-6}\left(\dfrac TK\right)^2$ ($T=$ temperature in Kelvin) then $\Delta G$ for given cell reaction at $27^\circ$C will be $\dfrac12$H$_2$ $+$ AgCl $\to$ Ag $+$ H$^+$ $+$ Cl$^-$
\begin{qoptions}
\opt{A}{$-0.44$ F}
\opt{B}{$-0.4$ F}
\opt{C}{$-0.22$ F}
\opt{D}{$-0.33$ F}
\end{qoptions}
\qn{2} For above cell reaction, entropy change per mole of Ag at $27^\circ$C will be (Use : $1$F $=96500$ C)
\begin{qoptions}
\opt{A}{$-115.8$ J/K}
\opt{B}{$-231.6$ J/K}
\opt{C}{$-57.6$ J/K}
\opt{D}{$-463.2$ J/K}
\end{qoptions}
\paraq{3}{4}
The hydrolysis reactions of halides can take place at room temperature or slightly higher temperature when the vacant orbitals is/are available on the central atom and some times on the bonded atoms.
\qn{3} Which compounds undergo partial hydrolysis :
\begin{qoptions}
\opt{A}{BF$_3$}
\opt{B}{SiCl$_4$}
\opt{C}{SiF$_4$}
\opt{D}{SOCl$_2$}
\end{qoptions}
\qn{4} Which of the following can be hydrolysed at room temperature :
\begin{qoptions}
\opt{A}{PCl$_3$}
\opt{B}{PCl$_5$}
\opt{C}{NF$_3$}
\opt{D}{NCl$_3$}
\end{qoptions}
\sect{SEC--3}
\qn{1} A gaseous mixture of H$_2$ and He in the mass ratio $1:4$ is successively effused through a series of containers connected through narrow tubes of negligible volume and fitted with valves. The valves are successively opened for a very small time interval, dt. Find the minimum number of total containers required to obtain a molar ratio $8:1$ for H$_2$ and He after the effusion process.
\qn{2} One mole of an diatomic ideal gas in an insulated cylinder fitted with a massless, frictionless movable piston is expanded gradually from $1$ L to $32$ L. If the initial temperature of gas is $1000$ K then find magnitude of change in Gibb's Free energy in kcal. (Molar entropy of gas $=3$ cal mol$^{-1}$/k$^{-1}$) (Use, $R=2$ cal mol$^{-1}$ k$^{-1}$) (Assume molar entropy of a gas is constant)
\qn{3} The activation energy of a reaction is given as $E_a\,(\mathrm{cal\,mol^{-1}})=0.01T^2+\dfrac{20}{7}T$ where $T$ is temperature in kelvin. Find value of $\ell n\dfrac{k_2}{k_1}$ where $k_1$ and $k_2$ are rate constants at temperature $200$ K and $400$ K respectively. (Use $\ell n\,2=0.7$ and $R=2$ cal mol$^{-1}$ k$^{-1}$)
\qn{4} A polyhydroxy compound having molecular mass $150$, on reacting the compound with CH$_3$COCl, the precipitate obtained has molecular mass of $318$. Calculate number of OH group present in compound :
\qn{5} Examine the given structure and select those that satisfy given conditions
\qfig{p2_chem_s3_q05_1}
Find out no. of compounds which give S$_\mathrm{N}$2 reaction with in CH$_3$S$^{\ominus}$Na$^{\oplus}$ in DMF faster with respect to CH$_3$CH$_2$Br
\qn{6} (a) Total number of stereoisomers of given compound are $x$
\qfig{p2_chem_s3_q06_1}
(b) Number of optically active stereoisomers of $1,2,3,4,5,6$-hexachlorocyclohexane is $y$. Then $x-y$ is
\qn{7} In the list of compounds given below, identify the number of compounds that have d$\pi-$p$\pi$ bonding in the molecular structure with sp$^2$ hybridization of the underlined atom.\\
\underline{C}Cl$_4$, \underline{P}OCl$_3$, \underline{S}O$_2$, \underline{S}O$_3$, H$_2$\underline{S}O$_4$, \underline{S}OCl$_2$, \underline{S}O$_2$Cl$_2$, \underline{B}F$_3$, \underline{B}Cl$_3$, \underline{N}(SiH$_3$)$_3$
\qn{8} How many of the following is/are yellow and water insoluble substances ?\\
Ag$_2$CrO$_4$, PbCrO$_4$, Na$_2$CrO$_4$, BaCrO$_4$, CrO$_2$Cl$_2$, (NH$_4$)$_2$Cr$_2$O$_7$, (NH$_4$)$_3$PO$_4$.12MoO$_3$, K$_3$[Co(NO$_2$)$_6$], AgI, PbI$_2$
\partsec{3}{MATHEMATICS}
\sect{SEC--1}
\qn{1} If
{\small\[f(x)=1-(\sin2x)\sqrt[4]{2\sqrt{\tan^2x\left(\tan^2x+6\right)+9}+\tan^4x-5},\]}
where $x\in\left(\dfrac{3\pi}2,\,2\pi\right]$
\begin{qoptions}
\opt{A}{$f(x)=1-2\sin x$}
\opt{B}{$f(x)=1+2\cos x$}
\opt{C}{$f(x)$ can take only $2$ integral value}
\opt{D}{$f(x)$ can take only $3$ integral value}
\end{qoptions}
\qn{2} $x+y=2$ and $x-y=2$ are tangents on a parabola at $(1,1)$ and $(4,2)$ respectively. Which of the following is/are correct
\begin{qoptions}
\opt{A}{Equation of directrix is $x+3y=2$}
\opt{B}{Equation of axis is $3x-y=5$}
\opt{C}{Focus of the parabola is at $\left(\dfrac85,\dfrac65\right)$}
\opt{D}{Vertex of the parabola is at $\left(\dfrac{33}{20},\dfrac{13}{20}\right)$}
\end{qoptions}
\qn{3} Probability of $n$ heads in $2n$ tosses of a fair coin can be given by
\begin{qoptions}
\opt{A}{$\displaystyle\prod_{r=1}^{n}\left(\dfrac{2r-1}{2r}\right)$}
\opt{B}{$\displaystyle\prod_{r=1}^{n}\left(\dfrac{n+r}{2r}\right)$}
\opt{C}{$\displaystyle\sum_{r=0}^{n}\left(\dfrac{{}^nC_r}{2^n}\right)^2$}
\opt{D}{$\dfrac{\displaystyle\sum_{r=0}^{n}\left({}^nC_r\right)^2}{\displaystyle\left(\sum_{r=0}^{2n}{}^{2n}C_r\right)}$}
\end{qoptions}
\qn{4} $P(n)$ denotes the number of natural number less than $n$, which are relatively prime with $n$, then
\begin{qoptions}
\opt{A}{$P\left(2^2.5\right)=2^3$}
\opt{B}{$P\left(2^3\right)=2^2$}
\opt{C}{$P\left(2^2.3\right)=2^2$}
\opt{D}{$P\left(2.3^2\right)=2^2.3$}
\end{qoptions}
\qn{5} If the line $y=2-x$ is tangent to the circle `$s$' at a point $P(1,1)$ and circle `$s$' is orthogonal to the circle $x^2+y^2+2x+2y-2=0$ then which of the following statement is true about circle `$s$'.
\begin{qoptions}
\opt{A}{The equation of the circle `$s$' is $x^2+y^2-x-y=0$}
\opt{B}{The length of the tangent from $(2,2)$ to the circle `$s$' is `$2$'}
\opt{C}{The equation of the director circle of circle `$s$' is $x^2+y^2-2x-2y+2=0$}
\opt{D}{The length of the common chord of given circle and circle `$s$' will be $\dfrac{\sqrt{10}}{3}$}
\end{qoptions}
\qn{6} If $I_m=\displaystyle\int_0^{m\pi}x\,|\sin x|\,e^{\cos4x}\,dx$ and $J_n=\displaystyle\int_0^{n\pi}|\cos x|\,e^{\cos4x}\,dx$, where $m,n\in I$ then (Note : $[k]$ denote GIF)
\begin{qoptions}
\opt{A}{$\left[\dfrac{I_1}{J_1}\right]=0$}
\opt{B}{$\left[\dfrac{I_2}{J_2}\right]=3$}
\opt{C}{$\left[\dfrac{I_2}{J_1}\right]=0$}
\opt{D}{$\left[\dfrac{I_3}{J_3}\right]=4$}
\end{qoptions}
\qn{7} Let $f(x)=\cos x\sin2x$, then
\begin{qoptions}
\opt{A}{$\min f(x)>-\dfrac79\ \forall x\in(-\pi,\pi)$}
\opt{B}{$\min f(x)>-\dfrac97\ \forall x\in(-\pi,\pi)$}
\opt{C}{$\min f(x)>-\dfrac19\ \forall x\in[-\pi,\pi]$}
\opt{D}{$\min f(x)>-\dfrac29\ \forall x\in[-\pi,\pi]$}
\end{qoptions}
\qn{8} Area bounded by $\sqrt{x+y+1}<\sqrt{y-x+2}$ \& $y<0$ is less than
\begin{qoptions}
\opt{A}{$\dfrac9{16}$}
\opt{B}{$\dfrac95$}
\opt{C}{$\dfrac54$}
\opt{D}{$\dfrac{16}5$}
\end{qoptions}
\sect{SEC--2}
\paraq{1}{2}
A square matrix $A=[a_{ij}]$ is called symmetric matrix if $a_{ij}=a_{ji}\ \ \forall\ i,j$ and a square matrix $A=[a_{ij}]$ is called skew-symmetric matrix if $a_{ij}=-a_{ji}\ \forall\ i,j$.
\qn{1} If $A$ and $B$ are symmetric matrices of order $3$ then, Which of the following statements are ALWAYS TRUE.\\
I) $A^{2n+1}$ is a skew-symmetric matrix for $n\in\mathbb N$\\
II) $AB+BA$ is a symmetric matrix,\\
III) $AB-BA$ is a skew-symmetric matrix\\
IV) $\mathrm{Det}(B)$ is non zero.
\begin{qoptions}
\opt{A}{I only}
\opt{B}{II \& III}
\opt{C}{I, II \& III}
\opt{D}{I, II, III \& IV}
\end{qoptions}
\qn{2} $A=\begin{bmatrix}0&2&-3\\-2&0&5\\3&-5&0\end{bmatrix}$, $X=\begin{bmatrix}x\\y\\z\end{bmatrix}$, $B=\begin{bmatrix}0\\0\\0\end{bmatrix}$ Then $AX=B$, will have
\begin{qoptions}
\opt{A}{Exactly one solution for $X$}
\opt{B}{No solution for $X$}
\opt{C}{Multiple solutions for $X$}
\opt{D}{None of these}
\end{qoptions}
\paraq{3}{4}
Linked comprehensive type passage: Equation of line $L=0$ is given by $2x-y+z-2=0=x+2y-z-3$, then
\qn{3} The equation of the plane containing the line $L=0$ and passing through the point $(3,2,1)$ is $ax+by+cz+1=0$, then $a-b+c$ is divisible by
\begin{qoptions}
\opt{A}{$3$}
\opt{B}{$4$}
\opt{C}{$2$}
\opt{D}{$6$}
\end{qoptions}
\qn{4} Vector equation of the plane which passes through the point $(-1,3,2)$ and is perpendicular to line $L=0$ is
\begin{qoptions}
\opt{A}{$\vec r\cdot\left(\hat i-3\hat j-5\hat k\right)+20=0$}
\opt{B}{$\vec r\cdot\left(\hat i-3\hat j+5\hat k\right)=0$}
\opt{C}{$\vec r\cdot\left(\hat i+3\hat j+5\hat k\right)-18=0$}
\opt{D}{$\vec r\cdot\left(-\hat i+3\hat j-5\hat k\right)+3=0$}
\end{qoptions}
\sect{SEC--3}
\qn{1} A tangent is drawn to the curve, $\dfrac{x^2}{16}+\dfrac{y^2}{9}=1$ at the point $P$ meeting the co-ordinate axis in $T$ and $t$. If perpendicular from origin $O$ to the about tangent meet at $Y$, then value of the product $(Tt)(PY)$ is equal to
\qn{2} A continuous function $f:\left[0,\dfrac34\right]\to\mathbb R$ satisfies the equation $f(x)=\left(1+x^2\right)\left(1+\displaystyle\int_0^x\dfrac{f^2(t)}{1+t^2}\,dt\right)$, then $\left[f\left(\dfrac12\right)\right]$ is

[Note : $[K]$ denotes greatest integer less than or equal to $K$.]
\qn{3} The quadratic polynomial $P(x)$ satisfy the equation $P(x)-P(x+1)=2x-1\ \forall\ x\in\mathbb R$. If the greatest value of $P(x)$ in $(-\infty,0]$ is $2$, then value of $P(-1)+P(1)$ is equal to
\qn{4}
{\small\begin{gather*}
S_n=\lim_{x\to\infty}\left(P^{1/n}-x\right),\ n\in\mathbb N,\ \text{where}\\
P=\left(x+\tfrac1{\sqrt{4n^2}}\right)\left(x+\tfrac1{\sqrt{4n^2-1^\alpha}}\right)\left(x+\tfrac1{\sqrt{4n^2-2^\alpha}}\right)\\
\cdots\left(x+\tfrac1{\sqrt{4n^2-(n-1)^\alpha}}\right)
\end{gather*}}
and for $\alpha=1$, $L=\displaystyle\lim_{n\to\infty}(n.S_n)$ and for $\alpha=2$, $M=\displaystyle\lim_{n\to\infty}(n.S_n)$ then value of $[L+M]$ is equal to

[Note : $[y]$ denotes the greatest integer function of $y$.]
\qn{5} Let a real valued function $f(x)=x^4+4x^3-8x^2+P$, $P\in\mathbb R$. If $\lambda$ is the least value of $P$ for which $g(x)=|f(x)|$ is differentiable $\forall\ x\in\mathbb R$, then value of $\dfrac\lambda{16}$ is equal to
\qn{6} Let $A$ and $B$ be two sets such that $n(A)=4$ and $n(B)=3$. If $(x,1)$, $(y,2)$, $(z,1)$ and $(2,3)$ are in $A\times B$, where $x,y,z$ are distinct elements. If a function $f:A\to B$ defined \& number of onto function from set $A$ to set $B$ is $\lambda$, then value of $\dfrac\lambda6$ is equal to
\qn{7} Let $3$ vector $\vec\ell,\vec m$ \& $\vec n$ are such that $\left|\vec\ell\right|=3$, $|\vec m|=\dfrac12$ \& $\vec\ell.\vec m=\vec m.\vec n=\vec n.\vec\ell=0$. If $\left[\vec\ell\times\left(\vec m\times\vec\ell\right)\quad\vec n\times\vec m\quad\vec n\right]=36$, then value of $|\vec n|$ is equal to
\qn{8} Let a circle with centre $O$ and radius $k$. $3$ points $P,Q,R$ are chosen on a circle such that if length of chord $PQ$, $QR$, $RP$ are $x$, $y$ \& $z$ respectively, $x^2+y^2+z^2=5$, $\angle QOR=60^\circ$ and Area of $\triangle PQR$ is $\dfrac{\sqrt3}4$ unit, then the value of radius of circle $k$ will be equal to
\end{multicols}
\end{document}
